% ============================================================
% oneline_mains.tex — mains distribution after the retrofit.
%
%   Motor branch (unchanged by the retrofit):
%     DISC -> 30 A -> M1 (Size 1) -> OL heater -> motor
%
%   Control branch (RETROFIT ADDITION):
%     control circuit -> VONVOFF receiver AC IN L/N
%                     -> AC OUT L into the seal-in rung
%
%   Supervisor supply (RETROFIT ADDITION, off-machine):
%     a separate 120 V wall outlet -> UL-listed USB charger
%                                  -> 5 V SELV -> ESP32
%
% The supervisor supply is deliberately NOT tapped from the
% machine. ARCHITECTURE.md's end-state design does tap it, behind
% a 250 mA fuse and an isolated PSU module, so that opening the
% machine disconnect de-energises the supervisor. That module
% (TASK-2) has not been bought, so today the ESP32 runs off a
% wall wart — which means the disconnect does NOT kill it. Called
% out in the legend, because it changes what "locked out" means.
%
% Regenerate with:  make oneline_mains.svg
% ============================================================
\def\pgfsysdriver{pgfsys-dvisvgm.def}
\documentclass[border=10pt,dvisvgm]{standalone}
% ============================================================
% tsstyle.tex — shared preamble for every sheet in this directory.
%
% Inputted by each sheet after \documentclass. Sheets stay pure
% geometry; everything about how the geometry *looks* lives here,
% so a style change is a one-file edit.
% ============================================================

\usepackage[T1]{fontenc}
\usepackage[utf8]{inputenc}
\usepackage[scaled=0.92]{helvet}
\renewcommand{\familydefault}{\sfdefault}

\usepackage[american,siunitx,RPvoltages]{circuitikz}
\usetikzlibrary{calc,positioning,fit,backgrounds,arrows.meta,decorations.pathreplacing}

% ------------------------------------------------------------
% Palette. Deliberately small: ink for conductors, mute for
% annotation, and exactly two semantic colours — mains (danger)
% and SELV (safe). Colour carries meaning on these sheets; it is
% never decoration.
% ------------------------------------------------------------
\definecolor{tsink}{HTML}{16181D}    % conductors, symbol bodies
\definecolor{tsmute}{HTML}{6E7178}   % secondary annotation
\definecolor{tsrule}{HTML}{C7CAD1}   % boxes, leader lines
\definecolor{tsmains}{HTML}{B3261E}  % mains potential — danger
\definecolor{tsselv}{HTML}{15607A}   % SELV / low-voltage rail
\definecolor{tssafe}{HTML}{1E6B45}   % fail-safe / passive layer
\definecolor{tspanel}{HTML}{F5F6F7}  % block fills
\definecolor{tspanelb}{HTML}{E4E7EB} % block fills, alternate

% ------------------------------------------------------------
% Global circuitikz sizing. Bipoles are scaled down from the
% default because these sheets carry a lot of elements per unit
% area; the default American resistor is comically large next to
% an IC block.
% ------------------------------------------------------------
\ctikzset{
  bipoles/length=0.85cm,
  resistors/scale=0.85,
  capacitors/scale=0.8,
  label/align=straight,
  font=\footnotesize,
}

\tikzset{
  % Conductors -------------------------------------------------
  wire/.style        = {tsink, line width=0.9pt, line cap=round},
  buswire/.style     = {tsink, line width=1.5pt, line cap=round},
  mainswire/.style   = {tsmains, line width=1.5pt, line cap=round},
  selvwire/.style    = {tsselv, line width=1.2pt, line cap=round},
  % Functional blocks -----------------------------------------
  block/.style       = {draw=tsink, line width=1pt, fill=tspanel,
                        rounded corners=1.5pt, align=center,
                        inner sep=6pt, font=\small\bfseries},
  subblock/.style    = {draw=tsrule, line width=0.7pt, fill=white,
                        rounded corners=1.5pt, align=center,
                        inner sep=4pt, font=\scriptsize},
  % Text ------------------------------------------------------
  pinlbl/.style      = {font=\scriptsize, tsink, inner sep=2pt},
  siglbl/.style      = {font=\scriptsize\itshape, tsmute, inner sep=1.5pt},
  netlbl/.style      = {font=\scriptsize\bfseries, tsink, inner sep=2pt},
  note/.style        = {font=\scriptsize, tsmute, align=left},
  danger/.style      = {font=\scriptsize\bfseries, tsmains, align=left},
  % Isolation boundary ----------------------------------------
  isoline/.style     = {tsmains, line width=1pt, dash pattern=on 5pt off 3pt},
  % A wire that crosses another without connecting. The white
  % under-stroke reads as an over-pass, so a crossing can never be
  % mistaken for a junction (junctions are always a filled dot).
  overpass/.style    = {preaction={draw, white, line width=3.4pt}},
  % Sheet furniture -------------------------------------------
  titleblock/.style  = {draw=tsink, line width=1pt, fill=white,
                        align=left, inner sep=7pt, font=\scriptsize},
  legendbox/.style   = {draw=tsrule, line width=0.7pt, fill=white,
                        align=left, inner sep=7pt, font=\scriptsize},
}

% ------------------------------------------------------------
% \pinrow{<x>}{<y>}{<dir>}{<label>}  — a pin stub on a block
% edge. dir is -1 (left edge, stub runs left) or +1 (right edge).
% ------------------------------------------------------------
\newcommand{\pinstub}[5]{%
  \draw[wire] (#1,#2) -- ++(#3*0.55,0);
  \node[pinlbl, anchor=#4] at (#1+#3*0.10,#2+0.14) {#5};
}

% ------------------------------------------------------------
% \titleblocknode{<x>}{<y>}{<sheet>}{<subtitle>}{<source>}
% Standard title block, anchored at its south east corner.
% ------------------------------------------------------------
\newcommand{\sheettitle}[5]{%
  \node[titleblock, anchor=south east] at (#1,#2) {%
    {\normalsize\bfseries #3}\\[2pt]
    {\color{tsmute}#4}\\[4pt]
    {\color{tsmute}Table-saw thermal protection retrofit\hfill}\\
    {\color{tsmute}Generated from \texttt{#5} — do not edit the SVG}%
  };
}


% CircuiTikZ's `elmech` is a node shape, not a bipole — `to[elmech]`
% compiles silently and draws a bare wire. Draw the machine directly.
\newcommand{\motorsym}[2]{%
  \draw[mainswire] (#1,#2) circle (0.62);
  \node[font=\normalsize\bfseries, tsink] at (#1,#2) {M};
}

\begin{document}
\begin{circuitikz}[line width=0.9pt]
% Elements on this sheet are sparse, so they can afford to be larger
% than the shared default.
\ctikzset{bipoles/length=1.1cm}

% ============================================================
% Provenance banner
% ============================================================
\node[font=\scriptsize\bfseries, tssafe, anchor=south west] at (0,2.6)
  {DRAWN FOR THE HARDWARE ON HAND. The motor branch is untouched by
   the retrofit.};

% ============================================================
% Retrofit highlight — control branch
% ============================================================
\begin{scope}[on background layer]
  \fill[tssafe!7, rounded corners=4pt] (3.6,-8.0) rectangle (14.9,-2.9);
  \draw[tssafe!35, line width=0.8pt, rounded corners=4pt]
        (3.6,-8.0) rectangle (14.9,-2.9);
\end{scope}
\node[font=\scriptsize\bfseries, tssafe, anchor=north west] at (3.75,-8.12)
  {ADDED BY THE RETROFIT --- control branch};

% ============================================================
% Utility source
% ============================================================
\node[font=\small\bfseries, anchor=south east] at (-0.1,0.55) {Utility};
\node[note, anchor=north east] at (-0.1,0.45) {240\,V, 1\,$\Phi$};
\draw[mainswire] (0,0) -- (2.4,0);

% ============================================================
% Motor branch — L1 across the top
% ============================================================
\draw[mainswire] (2.4,0) to[nos] (3.8,0);
\node[font=\small\bfseries, anchor=south] at (3.1,1.15) {DISC};
\node[note, anchor=north, align=center] at (3.1,1.05)
  {fused disconnect,\\lockable};

\draw[mainswire] (3.8,0) -- (6.6,0);
\draw (5.4,0) node[circ, color=tsmains]{};
\node[note, anchor=south, align=center, tssafe] at (5.4,0.55)
  {control-circuit tap\\\scriptsize load side of DISC,\\line side of M1};

\draw[mainswire] (6.6,0) to[fuse] (8.0,0);
\node[font=\small\bfseries, anchor=south] at (7.3,1.15) {30\,A};
\node[note, anchor=north] at (7.3,1.05) {motor branch};

\draw[mainswire] (8.0,0) -- (9.2,0);
\draw[mainswire] (9.2,0) to[nos] (10.6,0);
\node[font=\small\bfseries, anchor=south] at (9.9,1.15) {M1};
\node[note, anchor=north, align=center] at (9.9,1.05)
  {contactor,\\NEMA Size 1};
\node[note, anchor=north, align=center, tssafe] at (9.9,-0.75)
  {held closed by the\\coil circuit --- see\\\texttt{ladder\_coil\_circuit.svg}};

\draw[mainswire] (10.6,0) -- (11.8,0);
\draw[mainswire] (11.8,0) to[R] (13.2,0);
\node[font=\small\bfseries, anchor=south] at (12.5,1.15) {OL};
\node[note, anchor=north] at (12.5,1.05) {overload heater};

\draw[mainswire] (13.2,0) -- (15.08,0);
\motorsym{15.7}{0}
\node[note, anchor=west, align=left] at (16.55,1.15)
  {3\,HP, 14.4\,A FLA\\145TC TEFC\\Class F insulation};

% Return leg down to L2
\draw[mainswire] (16.32,0) -- (18.4,0) -- (18.4,-9.0);

% ============================================================
% L2 return
% ============================================================
\draw[mainswire] (0,0) -- (0,-9.0);
\draw[mainswire] (0,-9.0) -- (18.4,-9.0);
\node[font=\small\bfseries, anchor=north] at (9.2,-9.1) {L2};

% ============================================================
% Control branch — the RF receiver's supply and switched output
% ============================================================
\draw[mainswire] (5.4,0) -- (5.4,-4.2);
\draw[mainswire] (5.4,-4.2) -- (6.6,-4.2);

\draw[block, fill=white] (6.6,-6.4) rectangle (11.8,-3.4);
\node[font=\small\bfseries, tsink, anchor=north] at (9.2,-3.55)
  {VONVOFF RF receiver};
\node[note, anchor=north, align=center] at (9.2,-4.05)
  {\texttt{B07CTL3TG6}\\AC 100--240\,V, momentary mode};
\node[pinlbl, anchor=west] at (6.72,-4.2) {AC IN L};
\node[pinlbl, anchor=west] at (6.72,-5.6) {AC OUT N};
\node[pinlbl, anchor=east] at (11.68,-4.2) {AC OUT L};

\draw[mainswire] (11.8,-4.2) -- (13.6,-4.2);
\draw (13.6,-4.2) node[ocirc, color=tsmains]{};
\node[note, anchor=west, align=left, tsmains] at (13.8,-4.2)
  {$\rightarrow$ STOP, then the\\seal-in rung. Full ladder\\on
   \texttt{ladder\_coil\_circuit.svg}.};

% N bus back to L2 — on its own drop so it cannot be mistaken for
% a continuation of the AC IN L conductor at x = 5.4
\draw[mainswire] (6.6,-5.6) -- (4.2,-5.6);
\draw[mainswire] (4.2,-5.6) -- (4.2,-9.0);
\draw (4.2,-9.0) node[circ, color=tsmains]{};

\node[note, anchor=north west, align=left, tssafe] at (4.6,-6.6)
  {Tapped ahead of the seal-in so the receiver is powered and listening
   before START is\\
   pressed. Read the A202C's coil voltage off the coil label before
   landing this ---\\
   the receiver takes either 120\,V or 240\,V, but the tap has to match
   what is there.};

% ============================================================
% Supervisor supply — a SEPARATE circuit, not the machine's
% ============================================================
\begin{scope}[on background layer]
  \fill[tssafe!7, rounded corners=4pt] (0,-15.6) rectangle (11.6,-11.4);
  \draw[tssafe!35, line width=0.8pt, rounded corners=4pt]
        (0,-15.6) rectangle (11.6,-11.4);
\end{scope}
\node[font=\scriptsize\bfseries, tssafe, anchor=north west] at (0.15,-15.72)
  {ADDED BY THE RETROFIT --- supervisor supply, off-machine};

\node[font=\small\bfseries, anchor=south west] at (0.3,-12.0)
  {120\,V wall outlet};
\node[note, anchor=north west, align=left] at (0.3,-12.1)
  {a separate branch circuit,\\not this machine's};
\draw[mainswire] (0.4,-13.1) -- (3.4,-13.1);
\draw[mainswire] (0.4,-14.1) -- (3.4,-14.1);

\draw[block, fill=white] (3.4,-14.9) rectangle (7.8,-12.3);
\draw[isoline] (5.6,-14.9) -- (5.6,-12.3);
\node[font=\scriptsize\bfseries, tsmains, anchor=north] at (4.5,-12.45)
  {PRIMARY};
\node[font=\scriptsize\bfseries, tsselv, anchor=north] at (6.7,-12.45)
  {SECONDARY};
\node[note, anchor=north, tsmains] at (4.5,-13.5) {120\,V AC};
\node[note, anchor=north, tsselv]  at (6.7,-13.5) {5\,V DC};
\node[font=\small\bfseries, tsink, anchor=north, align=center] at (5.6,-15.0)
  {UL-listed USB charger};

\draw[selvwire] (7.8,-13.1) -- (10.4,-13.1);
\draw[selvwire] (7.8,-14.1) -- (10.4,-14.1);
\draw (10.4,-14.1) node[ocirc, color=tsselv]{};
\draw (10.4,-14.1) node[ocirc, color=tsselv]{};
\node[note, anchor=west, tsselv, align=left] at (8.0,-12.6)
  {\textbf{+5\,V (SELV)} $\rightarrow$ ESP32 USB-C};
\node[note, anchor=west, tsselv, align=left] at (8.0,-14.6)
  {0\,V $\rightarrow$ common GND bus};

% ============================================================
% Danger banner
% ============================================================
\node[danger, anchor=north west, text width=8.4cm] at (0,-9.8) {%
  Everything drawn in red is at mains potential. Only the SELV pair
  right of the charger's isolation barrier is safe to touch. Lock out
  the disconnect before working on any of it.};

% ============================================================
% Legends
% ============================================================
\node[legendbox, anchor=north west, text width=8.1cm] at (0,-17.2) {%
  \textbf{What changed vs.\ the OEM Powermatic}\\[4pt]
  \begin{tabular}{@{}p{0.35cm}@{}p{7.1cm}@{}}
    --- & Klixon BEC2921 \emph{removed} from the motor internal
          circuit (was a manual-reset thermal in series with the
          winding lead)\\[3pt]
    $+$ & VONVOFF RF receiver in the control circuit, at the head of
          the seal-in rung\\[3pt]
    $+$ & ESP32 supervisor and NTC probe, powered from a wall
          outlet\\[3pt]
    $=$ & Motor branch untouched --- same disconnect, same 30\,A
          branch, same contactor, same overload heaters\\
  \end{tabular}
};

\node[legendbox, anchor=north west, text width=8.9cm] at (9.0,-17.2) {%
  \textcolor{tsmains}{\textbf{The disconnect does not kill the
  supervisor}}\\[4pt]
  The ESP32 is fed from a wall outlet on a different circuit, so
  opening the machine disconnect leaves it powered and transmitting.
  That is harmless --- the receiver is fed from the machine's own
  control circuit and dies with it, so nothing can close --- but do
  not read a lit LED as "the machine is live," or a dark one as "the
  machine is dead."\\[5pt]
  ARCHITECTURE.md's end-state design removes this by tapping the
  supervisor from the machine behind a 250\,mA fuse and an isolated
  AC-DC module, so one disconnect kills everything. That module is
  TASK-2 and has not been bought.\\[5pt]
  \textbf{Isolation on this sheet}\\[3pt]
  \begin{tabular}{@{}p{2.7cm}@{\hspace{4pt}}p{5.4cm}@{}}
    USB charger  & reinforced, by virtue of the listing\\[2pt]
    Optocoupler  & ESP32 to fob --- see
                   \texttt{esp32\_supervisor.svg}\\[2pt]
    Receiver     & none. It is line-powered throughout.\\
  \end{tabular}
};

% ============================================================
% Title block
% ============================================================
\node[titleblock, anchor=north west, text width=6.6cm] at (12.4,-9.8) {%
  {\normalsize\bfseries Mains distribution --- one-line}\\[2pt]
  {\color{tsmute}Motor branch, control branch, supervisor supply}\\[5pt]
  {\color{tsmute}Table-saw thermal protection retrofit}\\
  {\color{tsmute}Coil circuit: \texttt{ladder\_coil\_circuit.svg}}\\
  {\color{tsmute}Source: \texttt{oneline\_mains.tex}}%
};

\end{circuitikz}
\end{document}
